\section*{Wstęp}\label{Wstęp}

% Wymagania:
% Ma za zadanie wprowadzenie i uzasadnienie problematyki podejmowanej w pracy magisterskiej, który w bardzo zwięzły sposób powinien opisywać co do tej pory wiadomo na dany temat, które elementy są ew. dyskusyjne i wymagają rozwiązania badawczego

Suplementy diety stanowią obecnie istotny element codziennych wyborów zdrowotnych i żywieniowych, zwłaszcza wśród osób młodych, w tym studentów \cite{narayanan2021students,arikawa2023dietary,abdul2020influencing,sundgot2022explanations,lieberman2015patterns,steele2005dietary}. Ich popularność wynika m.in. z powszechnej dostępności oraz powszechnego przekonania, że mogą wspierać zdrowie, poprawiać samopoczucie lub uzupełniać niedobory wynikające ze sposobu odżywiania. Jednocześnie suplementy diety, mimo że często są postrzegane przez konsumentów jako środki o działaniu zbliżonym do produktów leczniczych, zgodnie z definicją prawną stanowią środki spożywcze służące uzupełnieniu normalnej diety, a nie leczeniu czy zapobieganiu chorobom \cite{suplementy-a-leki}. Ta różnica ma istotne znaczenie praktyczne, ponieważ odmienne są zarówno wymagania dotyczące obrotu tymi produktami, jak i zasady ich stosowania oraz poziom nadzoru nad ich bezpieczeństwem. W populacji studenckiej problem suplementacji nabiera szczególnego znaczenia. Okres studiów wiąże się często z nieregularnym trybem życia, wzmożonym stresem, ograniczoną ilością snu, a nierzadko także z nieprawidłowymi nawykami żywieniowymi. Czynniki te mogą sprzyjać sięganiu po suplementy diety jako łatwo dostępne rozwiązanie mające poprawić kondycję fizyczną i psychiczną. Jednak częstotliwość ich stosowania nie zawsze idzie w parze z odpowiednią wiedzą na temat składu, wskazań, możliwych działań niepożądanych oraz zasad bezpiecznego stosowania. W praktyce może to prowadzić do przyjmowania preparatów bez konsultacji ze specjalistą, w dawkach niezgodnych z zaleceniami lub równolegle z innymi środkami o podobnym składzie, co zwiększa ryzyko błędów suplementacyjnych. W literaturze przedmiotu wskazuje się, że suplementy diety są często stosowane przez studentów, przy czym najczęściej wybierane są preparaty witaminowo-mineralne, witaminowe, kapsułki z olejami z ryb oraz suplementy białkowe \cite{kobayashi2017prevalence,lieberman2015patterns,radwan2019prevalence}. Jako główne motywy ich stosowania wymieniane są najczęściej troska o zdrowie, chęć poprawy odporności, wyglądu lub ogólnej kondycji organizmu, a także wpływ stresu i presji związany ze studiowaniem \cite{sundgot2022explanations,steele2005dietary,lawand2023assessment}. Istotną rolę w kształtowaniu postaw wobec suplementacji odgrywają również źródła informacji, wśród których najczęściej wybieranym jest internet \cite{kobayashi2017prevalence,chiba2020educational,freeman2015dietary}. Jednocześnie z badań wynika, że poziom wiedzy studentów na temat suplementów bywa niewystarczający, co może sprzyjać podejmowaniu decyzji zdrowotnych bez pełnej świadomości potencjalnych konsekwencji \cite{chiba2020educational,elsahoryi2023prevalence,homan2018dietary}. Pomimo rosnącego zainteresowania suplementacją wśród młodych dorosłych nadal istnieje potrzeba dokładniejszego rozpoznania sposobu ich stosowania w populacji studenckiej. Nadal nie wszystkie aspekty tego zjawiska są wystarczająco dobrze opisane, zwłaszcza w zakresie zgodności stosowania suplementów z zaleceniami bezpieczeństwa, roli wiedzy w podejmowaniu decyzji o suplementacji oraz źródeł, z których studenci czerpią informacje. W związku z tym zasadne wydaje się podjęcie badań, które pozwolą ocenić częstość stosowania suplementów diety, najczęściej wybierane preparaty, motywy ich przyjmowania oraz poziom wiedzy studentów na ten temat w jednym wspójnym badaniu. Celem niniejszej pracy jest ocena znajomości i częstotliwości stosowania suplementów diety wśród studentów, ze szczególnym uwzględnieniem najczęściej wybieranych preparatów, powodów ich stosowania oraz źródeł wiedzy. Podjęcie tej problematyki ma znaczenie nie tylko poznawcze, lecz także praktyczne, ponieważ może przyczynić się do lepszego zrozumienia zachowań zdrowotnych studentów oraz wskazać obszary wymagające rozwinięcia w zakresie edukacji jak i profilaktyki.

